\scanpage{45}
% archive sheet number: 44
\uline{Formes extérieures types}

\margin{\(\deleted{E_0=\operatorname{Ext}(G_0)}\quad H_0=\operatorname{ad}(G_0)\)}

\(G_0\) un groupe \(\tfrac12\) simple type (sur \(Z\)), \(\Gamma_0=\underline{\operatorname{Aut}}_Z(G_0)\), \deleted{\(\unsure{E=E/E^\circ|\Gamma=\Gamma/\Gamma^\circ}\)}, on a \deleted{\(\unsure{E\rtimes|E H}\)}\inserted{\(\Gamma_0\simeq E_0 H_0\)}

\(E_0\) sous-groupe \unsure{conservant|normalisant} \(T_0\subset B_0\)

\(S\) schéma de base

\(G\) forme de \(G_0\) sur \(S\), correspond : un principal \unsure{homogène|hom.} \deleted{\(\unsure{P_G|Q_G|P}\)}\inserted{\(\unsure{Q_0|P_0}\)} sous \(\Gamma_S\), \deleted{\illegible{1}} \deleted{\(\unsure{Q_G|P_G|P}\)}\inserted{\(\unsure{P|\rho}=P_G\)} le principal homogène associé de groupe \(E\), \deleted{qui} est \unsure{parfois|partie|noté} par abus de langage l'élément correspondant de \(H^1(S,E_0)\).

\inserted{\deleted{\unsure{plus précisément} \(\unsure{G_0^P|G^P}\)}}

\(G_1=G_0^P\) déduit de \(G_0\) en tordant par \(P\)\inserted{ [pour l'instant, \unsure{par hypoth.|par not.} \(P=P_G\)]} : une telle \deleted{groupe} forme \unsure{sera|ser.} dite \uline{forme extérieure type} associée à \(P\).

[On dit que \(G\) est une forme extérieure type s'il est \uline{isom.} à \(G_{1}^{\unsure{P|*|\rho}}\), \deleted{\unsure{tel}} \deleted{\unsure{si|et}} \unsure{nécessairemt|nécessairement} \(P\simeq P_G\)]

Si on part de \(P\), [sans \(G\)], \uline{la donnée d'un \(G\) et d'un} \uline{isom} \(\deleted{\unsure{G}}P\mathrel{\overset{\varphi}{\overset{\sim}{\longrightarrow}}}\deleted{\unsure{G}}P_G\) \textquotedblleft\uline{équivaut}\textquotedblright\ \uline{à la donnée d'un fibré principal homogène sous} \(H_1=\deleted{\unsure{G}}H_0^P\) ; la

% slanted
\margin{Distinguer entre les tores maximaux de \(G_0\) et \(H_0\), de \unsure{de|dans} \(G_1\) et \(H_1\), \ldots}

Classification des objets \(G\) à isom. près : [un isom de \((G,\varphi)\) et \((G',\varphi')\) étant donné : un isom \(\psi:G\to G'\) tel que l'isom correspondant \(P_G\xrightarrow{\sim}P_{G'}\) soit compatible avec les isom \(P_G\) et \(P_{G'}\) avec \(P\) \unsure{fixé|fix.} ; \uline{un tel} \deleted{\unsure{cet iso|cette}} \uline{isomorphisme est dit} \uline{intérieur}] \deleted{revient} \deleted{à \unsure{classifiant|classification}} est donnée par \(H^1(S,\deleted{\unsure{G}}H_{1}^{\unsure{P|1}})\). On a

% slanted
\margin{Définir ainsi la notion d'isom. \uline{\unsure{intérieur|extérieur}} d'un \unsure{gpe en lui-même|groupe en lui-même}.}

\[
H^1(S,\deleted{\unsure{G}}H_1^{\unsure{P|1}})\longrightarrow H^1(S,\Gamma)=\text{ens des classes \unsure{à isom|d'isom}}
\]

près de formes de \(G\) sur \(S\)

\(\text{\textquotedblright}\) \unsure{ens|ensemble|son} des classes : isom près de \textquotedblleft formes \deleted{relatives} intérieures relatives à \(P\)\textquotedblright
